% Scope: Explain the magnet, gradient and RF subsystems as physical implementations of earlier models; connect imperfections to encoding and safety limits.
% Status: architecture only; no chapter prose or completed problems yet.
% Prerequisites: Chapters 1, 4 through 6 and 14.
% Planned worked examples: Gradient slew, receive-noise budget and digitization timing.
% Planned figures: Magnet cross-section, Maxwell/Golay coils and receiver chain.
% Planned challenge problems: Trace a hardware fault through k-space; allocate a receiver noise budget.
\chapter{MRI Hardware}
\label{ch:17}

\section{Main magnet and field control}
\label{sec:17:main-magnet-and-field-control}
\subsection{Superconductivity, persistent currents and cryogenics}
\subsection{Homogeneity, drift and passive shimming}
\subsection{Active shims and field monitoring}

\section{Gradient subsystem}
\label{sec:17:gradient-subsystem}
\subsection{Maxwell pairs, Golay coils and linearity}
\subsection{Amplifiers, inductance, amplitude and slew}
\subsection{Eddy currents, mechanical vibration and PNS links}

\section{RF transmit subsystem}
\label{sec:17:rf-transmit-subsystem}
\subsection{Birdcage coils, quadrature and loading}
\subsection{B1-plus fields and transmit calibration}
\subsection{Parallel transmission and coupling}

\section{RF receive subsystem}
\label{sec:17:rf-receive-subsystem}
\subsection{Surface coils and phased arrays}
\subsection{Coil sensitivity, sample noise and preamplifiers}
\subsection{Decoupling, matching and transmit-receive protection}

\section{Receiver and control chain}
\label{sec:17:receiver-and-control-chain}
\subsection{Mixing, filtering and complex demodulation}
\subsection{ADC bandwidth, dynamic range and quantization}
\subsection{Timing, synchronization and gradient monitoring}

\section{System-level limitations}
\label{sec:17:system-level-limitations}
\subsection{Field strength and noise regimes}
\subsection{Calibration drift and quality assurance}
\subsection{Hardware constraints in sequence design}

\section{Worked examples}
\label{sec:17:worked-examples}
% Planned calculations: Gradient slew, receive-noise budget and digitization timing.
\subsection{Derivation and quantitative calculation}
\subsection{Assumptions, limiting cases and scanner interpretation}

\section{Challenge problems}
\label{sec:17:challenge-problems}
% Problem design: Trace a hardware fault through k-space; allocate a receiver noise budget.
% Use challengeproblem and stable labels prob:17:<topic>.
% Complete solutions belong in Appendix E, section for this chapter.
% Use \begin{solution}{prob:17:<topic>} ... \end{solution} there.
\subsection{Derivation and quantitative design}
\subsection{Model criticism and integrated reasoning}
